{\mfontCalibri
Since this is not a typical Gita study book, the following approach will help maximise its benefits.
\begin{description}
  \item[\mfontCalibri{Sequential Reading of \emph{Shlokas}} \textemdash] This book provides all the Gita \emph{shlokas} (verses) along with their meanings. Read through the verses in the given order before engaging with the \emph{Mananam} (Contemplation) and \emph{inspiration} sections.
  \item[\mfontCalibri{Contextual Understanding} \textemdash] While the \emph{Mananam} questions focus on specific \emph{shlokas}, it is beneficial to read the preceding verses to grasp the broader context.
  \item[\mfontCalibri{Background in Gita \emph{Shravanam}} \textemdash] Ideally, one should have studied the Gita through structured exposition by a traditional Vedanta teacher (\emph{shravanam}). However, such \emph{shravanam} and prior formal study is not mandatory for using this book.
  \item[\mfontCalibri{Deep Reflection on \emph{Mananam} Questions} \textemdash] After reading the \emph{shlokas}, engage with the \emph{Mananam} questions thoughtfully. Read them multiple times, contemplating how they apply to your life and how their wisdom can be practiced or implemented.
  \item[\mfontCalibri{Writing and Revisiting}  \textemdash] Ponder upon the questions, and if necessary, revisit the \emph{shlokas}. Make notes in the space provided in the book or in a separate notebook, recording insights and reflections.
  \item[\mfontCalibri{Cultivating \emph{Jignāsu}}  \textemdash] If additional questions arise during contemplation, write them down as well. This enquiry (\emph{jignāsu}) will naturally fuel the quest for deeper understanding.
  \item[\mfontCalibri{Gaining Perspective}  \textemdash] Conclude by reading the \emph{inspiration} notes. Compare them with your interpretations—do they align with your insights or offer new perspectives?
  \item[\mfontCalibri{Informal Contemplation}  \textemdash] Beyond structured study, integrate contemplation into daily life. Select \emph{Mananam} questions that resonate with you, and reflect on them during free moments throughout the day.
  \item[\mfontCalibri{Prioritising Depth Over Quantity} \textemdash] The impact of this book lies not in how much is read, but in the quality and depth of contemplation. If you read for five minutes, dedicate at least 10–15 minutes to formal contemplation, allowing the teachings to permeate your thoughts informally throughout the day.
  \item[\mfontCalibri{Living the Teachings} \textemdash] Let the fragrance of the Gita infuse every aspect of your life, enriching your interactions and radiating wisdom to those around you.
\end{description} }
